\documentclass[conference]{IEEEtran}
\IEEEoverridecommandlockouts

\usepackage{cite}
\usepackage{amsmath,amssymb,amsfonts}
\usepackage{graphicx}
\usepackage{textcomp}
\usepackage{xcolor}
\usepackage{booktabs}
\usepackage{multirow}
\usepackage[T1]{fontenc}
\usepackage[hidelinks]{hyperref}

% Values marked "% TODO" are from one training seed (seed 0) or not measured yet.
% Replace them with the 3-seed results before submission (see outline_v2.md).

\begin{document}

\title{MA-X3D: A Lightweight Motion-Attentive 3D CNN\\
for Violence Detection in Surveillance Video}

\author{\IEEEauthorblockN{Nguyen Huu Hoang}
\IEEEauthorblockA{\textit{School of Information and Communications Technology} \\
\textit{Hanoi University of Science and Technology}\\
Hanoi, Vietnam \\
hoang.nh225972@sis.hust.edu.vn}
\and
\IEEEauthorblockN{Ngo Thanh Trung}
\IEEEauthorblockA{\textit{School of Information and Communications Technology} \\
\textit{Hanoi University of Science and Technology}\\
Hanoi, Vietnam}
}

\maketitle

\begin{abstract}
Automatic violence detection is useful only if it can run on every stream of
a surveillance installation, yet the most accurate models on the RWF-2000
benchmark have tens to hundreds of millions of parameters. We ask how much
of this accuracy a compact 3D convolutional network can retain within a
budget of a few GFLOPs per clip. The proposed MA-X3D adds two modules to
X3D-M, both initialised as the identity: a Motion Attention gate that
rescales mid-level features according to a multi-scale frame-difference
map, and a reparameterised $5\times5$ channel-wise kernel in the two
earliest residual stages, whose outer ring is trained in its own
learning-rate group and fused into a single kernel for inference. Together
they add 36K parameters. The network is trained with discriminative
fine-tuning and with distillation from a VideoMAE-B teacher that receives
the same augmented clips. Because RWF-2000 has no validation split and its
training split contains copies of test clips, we remove duplicates and take
every design decision by grouped four-fold cross-validation, so that the
test split is used once. Under this protocol a single MA-X3D model with
3.0M parameters and 5.5 GFLOPs reaches 87.6\% accuracy and 0.875 F1 on the
test split, compared with 86.1\% for conventional fine-tuning of the same
network and 91.2\% for the teacher, which is 29 times larger.
% TODO: one sentence on Tables IV-V (modules vs controls) once they are complete.
\end{abstract}

\begin{IEEEkeywords}
violence detection, video action recognition, surveillance video,
efficient 3D CNN, motion attention, knowledge distillation
\end{IEEEkeywords}

%======================================================================
\section{Introduction}
\label{sec:intro}

Closed-circuit television is now installed at a scale that human
monitoring cannot match. An estimated 770 million surveillance cameras were
in operation by the end of 2019~\cite{cnbc2019}, while the ability of an
observer to detect rare signals during continuous monitoring starts to
decline within the first twenty minutes~\cite{mackworth1948}. Automatic
violence detection addresses this mismatch by flagging the few streams that
need attention and leaving the decision to the operator.

The task is harder than generic action recognition in three respects.
Violent events are brief and involve fast, irregular movement; the people
involved may occupy a small part of a wide-angle frame; and many
non-violent activities, such as contact sports or play, produce similar
motion. Methods have followed the development of action recognition, from
hand-crafted motion descriptors~\cite{idt,vif,hockey} to two-stream
networks with optical flow~\cite{twostream}, 3D convolutional
networks~\cite{c3d,i3d,slowfast,x3d} and video
transformers~\cite{timesformer,videoswin}.

A practical detector must also be cheap, because it runs continuously on
every camera of an installation. Recent accuracy gains on the RWF-2000
benchmark~\cite{rwf2000} have come from models far outside this budget.
CUE-Net~\cite{cuenet} reports 94.0\% with 354M parameters and about
5.8~TFLOPs per clip, and masked-autoencoder transformers report 91.8--92.4\%
with 67M to 305M parameters~\cite{videomae,mixvit}. Lightweight
detectors~\cite{rwf2000,sepconvlstm,mdcn,biltma} report 87--90\% with a few
million parameters or fewer. This paper works in the second regime and asks
what limits a compact 3D convolutional network in it.

We see three limitations. First, a 3D backbone receives motion only
implicitly, as changes of appearance between frames that it must learn to
extract. The frame difference is a direct and nearly free motion cue that
is effective on this benchmark~\cite{sepconvlstm}, but prior detectors
process it in a second stream, which doubles the backbone cost. Second, the
receptive field of the early stages is narrow. X3D uses $3\times3$ spatial
kernels in all its channel-wise convolutions, including the two stages that
run at $56\times56$ and $28\times28$, where small and fast displacements are
still resolved. Third, with 1{,}600 training clips, fine-tuning a small
pre-trained backbone falls well short of large pre-trained transformers,
and the size of this gap is hard to judge from published numbers, because
the benchmark has no validation split and the procedure used to select
checkpoints is often not reported.

We address the first two limitations with two modules attached to
X3D-M~\cite{x3d}. Motion Attention converts a multi-scale frame-difference
map into four learnable temporal modes and uses them to rescale the Res4
features per channel and per position, within a bounded range. The
wide-kernel convolution enlarges the channel-wise kernels of Res2 and Res3
from $3\times3$ to $5\times5$ by adding a zero-initialised outer ring. The
ring is kept as a separate tensor during training, so that it can use a
larger learning rate than the pre-trained centre, and is summed into one
kernel afterwards. Both modules are initialised as the identity, so
training starts from the pre-trained function, which matters when 1{,}600
clips have to adapt a network pre-trained on Kinetics. For the third
limitation we change the training procedure rather than the network: a
higher learning rate for the pre-trained layers, and long
distillation~\cite{hinton2015,beyer2022} from a VideoMAE-B teacher that
receives the same augmented clip as the student. The teacher is used only
during training.

To compare these choices reliably we also revise the evaluation. Nine
training clips of RWF-2000 are copies of test clips. We remove them,
together with duplicate pairs inside the training split, and take every
decision (architecture, training recipe, number of views, threshold) on
the out-of-fold predictions of a grouped four-fold cross-validation, so
that the test split is scored once. Under this protocol a single MA-X3D
model reaches 87.6\% test accuracy with 3.0M parameters and 5.5 GFLOPs,
1.5 points above conventional fine-tuning of the same network and 3.6
points below the teacher, which has 29 times more parameters. Most of this
gain comes from the higher learning rate; distillation raises out-of-fold
accuracy by a further 1.2 points, but its effect on the test split is
within the variation between models. The learned outer ring carries 30\%
of the kernel magnitude, and removing it at test time lowers accuracy to
80.6\%.

Our contributions are the following.
\begin{itemize}
  \item A Motion Attention module that injects a multi-scale
        frame-difference cue into a single 3D backbone through learnable
        temporal modes and a bounded, identity-initialised gate, at 2.5K
        parameters and 0.02 GFLOPs.
  \item A wide-kernel reparameterisation of the channel-wise convolutions
        of the two earliest stages, trained in its own learning-rate group
        and fused exactly into a single $5\times5$ kernel for inference
        (34K parameters, 0.72 GFLOPs).
  \item A training recipe for a compact video backbone on a small
        dataset, a higher backbone learning rate and long, view-consistent
        distillation from a video transformer, which raises single-model
        test accuracy from 86.1\% to 87.6\% without changing the deployed
        network.
  \item A leakage-free protocol for RWF-2000, with duplicate removal,
        grouped cross-validation and a single use of the test split. The
        split files and code are released.
\end{itemize}

Section~\ref{sec:related} reviews related work, Section~\ref{sec:method}
presents MA-X3D and its training, Section~\ref{sec:protocol} describes the
protocol, Section~\ref{sec:results} reports the results, and
Section~\ref{sec:conclusion} concludes.

%======================================================================
\section{Related Work}
\label{sec:related}

\subsection{Violence detection}

Hand-crafted methods encode motion along dense
trajectories~\cite{idt}, summarise changes of optical-flow
magnitude~\cite{vif}, or use acceleration descriptors~\cite{hockey}. They
are cheap but sensitive to viewpoint and lighting. Two-stream
networks~\cite{twostream} learn from RGB and optical flow; the flow is often
more expensive than the network. 3D convolutional
networks~\cite{c3d,i3d,slowfast,x3d} learn spatio-temporal features from
RGB alone, and video transformers~\cite{timesformer,videoswin} extend
self-attention to space and time at the price of large-scale pre-training.

On RWF-2000, ConvLSTM~\cite{convlstm} reports 77.00\% and
C3D~\cite{c3d} 82.75\%. The Flow Gated Network~\cite{rwf2000} gates an RGB
branch with an optical-flow branch and reaches 87.25\%.
SepConvLSTM~\cite{sepconvlstm} processes background-suppressed frames and
frame differences with two MobileNetV2 backbones and reaches 89.75\%. A
two-stream multi-dimensional network~\cite{mdcn} reaches 89.70\% below one
million parameters, and an EfficientNet-B0 with bidirectional motion
attention and temporal shift~\cite{biltma} reaches 90.25\% at 4.48M
parameters. VideoMAE transformers reach 91.8--92.4\%~\cite{videomae,mixvit}
and CUE-Net~\cite{cuenet} 94.00\%. X3D itself has been reported at 84.75\%~\cite{cuenet} and,
with person tubes and a different test split, at 84.02\%~\cite{pathak2024}. Most of these works report accuracy on
the official test split and do not state how the checkpoint was chosen; we
return to this point in Section~\ref{sec:protocol}.

\subsection{Efficient video backbones}

X3D~\cite{x3d} expands a small 2D network one axis at a time (duration,
frame rate, resolution, depth, width, bottleneck width) and keeps the
expansion with the best accuracy-to-cost ratio. X3D-M reaches 76.0\% top-1
on Kinetics-400 with 3.8M parameters and 6.2 GFLOPs per view, comparable to
SlowFast R50 with nine times fewer parameters. X3D-S and X3D-M share the
same parameter count and differ in input resolution, which makes X3D-M the
most accurate member of the family within the smallest parameter budget.

\subsection{Motion cues, gating and reparameterisation}

Under a static camera the absolute difference between frames isolates the
movement of people, at a fraction of the cost of optical flow.
Multiplicative gating~\cite{se,cbam} re-weights features with a learned
signal; when the last layer of the gate starts near zero, the module starts
as an identity and does not disturb pre-trained features. Structural
reparameterisation~\cite{repvgg,acnet} trains extra branches and fuses them
into one kernel for inference. Both cited methods fuse back into the
original kernel shape; our ring fuses into a larger kernel, so the added
spatial support remains.

\subsection{Knowledge distillation}

Distillation trains a small student to match the softened outputs of a
larger teacher~\cite{hinton2015}. Beyer~et~al.~\cite{beyer2022} show that
the student benefits when teacher and student see the same augmented view
and when the schedule is much longer than for standard training. We follow
both observations with a video transformer as teacher and a 3D CNN as
student, and train one teacher per cross-validation fold so that no
validation clip is seen by either model.

%======================================================================
\section{Method}
\label{sec:method}

\subsection{Problem formulation and overview}

Given a trimmed clip $X \in [0,1]^{3 \times T \times H \times W}$ of RGB
frames, with $T=16$ and $H=W=224$, the task is to predict whether it
contains violence, $y \in \{0,1\}$. The model $f_\theta$ returns two logits,
and the clip is labelled violent when the softmax probability of that class
is at least 0.5.

MA-X3D keeps the X3D-M backbone~\cite{x3d} unchanged and adds two modules
(Fig.~\ref{fig:pipeline}, Table~\ref{tab:arch}). The stem applies a
$1\times3\times3$ convolution with spatial stride 2 and a $5\times1\times1$
channel-wise temporal convolution. Four residual stages follow; each block
is a bottleneck made of a $1\times1\times1$ expansion to 2.25 times the
stage width, a $3\times3\times3$ channel-wise convolution (with spatial
stride 2 in the first block of a stage), squeeze-and-excitation~\cite{se}
and Swish, and a $1\times1\times1$ projection. The temporal resolution stays
at 16 throughout. The head projects to 432 channels, pools globally, maps
to 2048 features, and ends in a new two-way linear layer. The backbone is
pre-trained on Kinetics-400~\cite{kinetics}.

The design follows four principles. \textbf{P1 Lightweight}: no second
stream and no change to the backbone topology; the two modules add 36.2K
parameters (1.2\%) at inference. \textbf{P2 Identity initialisation}: at
the start of training the network computes the same function as the
pre-trained backbone, so the new parameters are learned as departures from
it. \textbf{P3 No inference overhead beyond the widened kernels}:
training-time structure is fused away. \textbf{P4 Training-only teacher}:
the distillation teacher is not part of the deployed model.

\begin{table}[t]
  \centering
  \caption{MA-X3D for a $16\times224\times224$ clip. Width: output channels
  (inner width of the bottleneck in brackets).}
  \label{tab:arch}
  \setlength{\tabcolsep}{3pt}
  \begin{tabular}{lcccl}
    \toprule
    Stage & Output $T{\times}H{\times}W$ & Width & Blocks & Change \\
    \midrule
    Input & $16\times224\times224$ & 3   & -- & motion map $m$ \\
    Stem  & $16\times112\times112$ & 24  & -- & -- \\
    Res2  & $16\times56\times56$   & 24 (54)  & 3  & $5\times5$ kernels \\
    Res3  & $16\times28\times28$   & 48 (108) & 5  & $5\times5$ kernels \\
    Res4  & $16\times14\times14$   & 96 (216) & 11 & Motion Attention \\
    Res5  & $16\times7\times7$     & 192 (432) & 7 & -- \\
    Head  & $1\times1\times1$      & 2   & -- & new classifier \\
    \bottomrule
  \end{tabular}
\end{table}

\begin{figure}[t]
  \centering
  % \includegraphics[width=\linewidth]{fig/pipeline.pdf}
  \caption{MA-X3D. The motion map is computed from the input frames and
  gates the Res4 features; Res2 and Res3 use $5\times5$ channel-wise
  kernels. During training (dashed) a VideoMAE-B teacher receives the same
  augmented clip and provides soft targets; it is not used at inference.}
  \label{fig:pipeline}
\end{figure}

\subsection{Input pipeline}

\textbf{Person-centred crop.} Offline, YOLOv8n~\cite{yolov8} detects people
on 12 of 64 frames spread over the clip, the centres of the detected boxes
are clustered with DBSCAN~\cite{dbscan} (radius $0.15\max(H,W)$, two points
minimum), and the box of the most populated cluster, enlarged by 20\% on
each side, is applied to all frames; the full frame is kept when no person
is found. Clustering rather than taking the union of all boxes keeps
bystanders from widening the crop. The crop is static: under a fixed camera
the background remains registered across frames, so the frame difference
reflects the movement of people rather than of the crop window. The frames
are denoised with a bilateral filter~\cite{bilateral}, contrast-enhanced
with CLAHE~\cite{clahe} on the luminance channel, and resized to
$224\times224$; 64 frames are stored per clip.

\textbf{Temporal sampling.} At test time the 16 frames are spread evenly
over the 64 stored frames. During training they are spread over a random
window covering 60--100\% of the clip, with a jitter of one frame, so that
training and testing see motion at the same time scale.

\textbf{Augmentation.} Every random parameter is drawn once per clip and
applied to all its frames, so that the frames stay registered and the
motion map contains only real motion: brightness ($\times0.5$--$1.5$),
saturation ($\pm30\%$) and value shifts, horizontal flip, random crop
(70--100\% of the side, $p=0.8$), rotation ($\pm10^\circ$, $p=0.8$), a
change of playback speed ($\times0.7$--$1.3$, $p=0.5$) and Gaussian blur
($p=0.2$). CutMix~\cite{cutmix} is applied to half of the batches with a
box that is the same in all frames of a clip.

\subsection{Motion Attention}
\label{sec:ma}

\textbf{Motion map.} Let $x_{c,t}$ denote channel $c$ of frame $t$ of the
input clip in $[0,1]$. The channel-averaged absolute difference at step
$s$ is
\begin{equation}
  d^{(s)}_t = \frac{1}{3} \sum_{c=1}^{3} \left| x_{c,t} - x_{c,t-s} \right|,
  \qquad s \in \{1, 2\},
  \label{eq:diff}
\end{equation}
with $d^{(s)}_t = d^{(s)}_s$ for $t < s$, so that the map keeps all $T$
frames. The motion map combines the two steps and saturates at $\delta$:
\begin{equation}
  m_t = \frac{1}{\delta}\min\!\left( \max\!\left(d^{(1)}_t,\, d^{(2)}_t\right),\; \delta \right),
  \qquad \delta = 0.2 .
  \label{eq:motion}
\end{equation}
The step-1 and step-2 differences respond to fast and to slower movement;
taking their maximum marks a location as moving if either responds and
keeps the map to one channel. Saturation limits the influence of flicker
and compression artefacts and maps typical differences to $[0,1]$. The map
is computed from the input frames before Kinetics normalisation, so it
measures intensity change independently of the normalisation statistics.

\textbf{Temporal modes.} A temporal convolution $w \in \mathbb{R}^{N \times
1 \times 3 \times 1 \times 1}$ maps $m$ to $N=4$ channels, $M = w * m$.
Because its spatial extent is one, each output channel is a temporal filter
of the motion signal and can represent, for instance, its onset or its
persistence. The first filter is initialised as the impulse $(0,1,0)$,
which reproduces $m$, and the others from $\mathcal{N}(0, 0.05^2)$.

\textbf{Gate.} The modes are resampled to the Res4 resolution by trilinear
interpolation, $\tilde{M} = \mathcal{R}(M) \in \mathbb{R}^{4 \times 16
\times 14 \times 14}$, and mapped to one gain per channel and position:
\begin{equation}
  G = \tanh\!\left( W_2\, \mathrm{ReLU}\!\left(W_1 \tilde{M}\right) + b_2 \right),
  \label{eq:gate}
\end{equation}
where $W_1 \in \mathbb{R}^{24 \times 4}$ and $W_2 \in \mathbb{R}^{96 \times
24}$ act as $1\times1\times1$ convolutions. The Res4 output $F \in
\mathbb{R}^{96 \times 16 \times 14 \times 14}$ is modulated as
\begin{equation}
  Y = F \odot (1 + G) .
  \label{eq:modulate}
\end{equation}

\textbf{Properties.} (i) \emph{Identity at initialisation}: $W_2$ is drawn
from $\mathcal{N}(0, 0.01^2)$ and $b_2 = 0$, so $G \approx 0$ and
$Y \approx F$ (P2). (ii) \emph{Bounded modulation}: $1 + G \in (0, 2)$, so
the motion cue can attenuate a feature or at most double it, but cannot
create a response where the backbone has none; appearance information is
re-weighted, not replaced. A multiplicative form is preferred to an
additive one for this reason. (iii) \emph{Joint gain}: $G$ varies over
channels, time and space, whereas squeeze-and-excitation~\cite{se} varies
over channels only and the spatial branch of CBAM~\cite{cbam} over
positions only.

\textbf{Placement.} The gate is applied after Res4, whose $14\times14$
features are semantic but still resolve where movement occurs. Res5, at
$7\times7$, is too coarse to localise motion, and Res2 and Res3 still encode
low-level texture. The module has 2{,}508 parameters and adds 0.017 GFLOPs
(Table~\ref{tab:complexity}).

\begin{figure}[t]
  \centering
  % \includegraphics[width=\linewidth]{fig/ma.pdf}
  \caption{Motion Attention. The multi-scale motion map is filtered into
  four temporal modes, resampled to the Res4 resolution, and turned into a
  gain $1+G \in (0,2)$ that multiplies the Res4 features.}
  \label{fig:ma}
\end{figure}

\subsection{Wide-kernel reparameterisation}
\label{sec:wk}

\textbf{Formulation.} Let $W_{\text{base}} \in \mathbb{R}^{C \times 1
\times 3 \times 3 \times 3}$ be the pre-trained channel-wise kernel of a
block in Res2 ($C=54$) or Res3 ($C=108$), with temporal, vertical and
horizontal extent 3. We replace it by
\begin{equation}
  W_{\text{eff}} = \mathcal{P}(W_{\text{base}}) + M_{\text{ring}} \odot \Delta ,
  \label{eq:rep}
\end{equation}
where $\mathcal{P}$ pads the two spatial dimensions with zeros to
$5\times5$, $\Delta \in \mathbb{R}^{C \times 1 \times 3 \times 5 \times 5}$
is initialised to zero, and $M_{\text{ring}}$ is a fixed binary mask equal
to one on the 16 outer positions of each $5\times5$ slice and zero on the
central $3\times3$. The spatial padding of the convolution grows by one, so
the output size is unchanged. At initialisation $W_{\text{eff}}$ equals the
pre-trained kernel (P2), and throughout training the mask keeps the centre
and the ring as disjoint sets of weights.

\textbf{Exact fusion.} Convolution is linear in its kernel,
$W_{\text{eff}} * X = \mathcal{P}(W_{\text{base}}) * X + (M_{\text{ring}}
\odot \Delta) * X$. After training, $W_{\text{eff}}$ is evaluated once and
stored as an ordinary $3\times5\times5$ channel-wise kernel, so inference
runs a single convolution per block (P3). In our runs the fused and the
unfused networks differ by less than $4\times10^{-7}$ in output
probability.

\textbf{Why a separate tensor.} The simpler alternative inflates the kernel
to a single dense $5\times5$ tensor whose ring starts at zero. The ring then
belongs to a pre-trained layer and receives the backbone learning rate.
Adam-type optimisers~\cite{adamw} divide each gradient by a running
estimate of its magnitude, so the size of an update is set by the learning
rate rather than by the gradient; at the backbone rate the ring can move no
further than the pre-trained weights around it. Keeping $\Delta$ as its own
tensor places it in the group of new parameters, with a learning rate ten
times that of the backbone, while the centre is fine-tuned at the backbone
rate. Section~\ref{sec:results} compares the two variants.

\textbf{Receptive field and cost.} A $k\times k$ spatial kernel widens the
receptive field by $k-1$ pixels of its feature map, so the change from
$3\times3$ to $5\times5$ doubles the contribution of each block. Over a
stage of $L$ blocks the widening grows from $2L$ to $4L$ pixels, from 6 to
12 in Res2 ($L=3$) and from 10 to 20 in Res3 ($L=5$). The modification
covers all eight blocks of the two stages and adds $48 \times (3 \cdot 54 +
5 \cdot 108) = 33{,}696$ parameters after fusion and 0.715 GFLOPs, most of
the added cost, because these stages run at the highest resolutions.
Widening Res4 and Res5 would cost more parameters, as their channels are
wider, while their receptive fields are already large relative to their
$14\times14$ and $7\times7$ maps.

\begin{table}[t]
  \centering
  \caption{Cost of each module at inference ($16\times224\times224$, ring
  fused; GFLOPs counted as multiply--accumulates with fvcore).}
  \label{tab:complexity}
  \setlength{\tabcolsep}{4pt}
  \begin{tabular}{lrr}
    \toprule
    Model & Parameters & GFLOPs \\
    \midrule
    X3D-M (two-way head)   & 2{,}978{,}772 & 4.732 \\
    + Motion Attention     & +2{,}508      & +0.017 \\
    + wide kernels         & +33{,}696     & +0.715 \\
    MA-X3D                 & 3{,}014{,}976 & 5.464 \\
    \bottomrule
  \end{tabular}
\end{table}

\begin{figure}[t]
  \centering
  % \includegraphics[width=\linewidth]{fig/wk.pdf}
  \caption{Wide-kernel reparameterisation. The pre-trained $3\times3$
  centre is padded to $5\times5$ and a zero-initialised ring $\Delta$,
  restricted by a fixed mask, is added and trained at the learning rate of
  the new parameters. After training the two are summed into one kernel.}
  \label{fig:wk}
\end{figure}

\subsection{Training}
\label{sec:training}

\textbf{Schedule and parameter groups.} Training has two phases
(Table~\ref{tab:phases}). In the probe phase the backbone is frozen and
only the new classifier and Motion Attention are trained, so that the new
components adapt to the pre-trained features before any pre-trained weight
changes. In the fine-tuning phase Res2, Res3, Res5 and the head are
trained at the backbone rate, and the new parameters, including the rings
$\Delta$, at a rate ten times higher. The stem and Res4 stay frozen, and so
do the BatchNorm statistics of all frozen layers. The optimiser is rebuilt
at the phase change so that the moment estimates start afresh.

\begin{table}[t]
  \centering
  \caption{Trainable parameters and learning rates. The stem and Res4 are
  frozen in both phases.}
  \label{tab:phases}
  \setlength{\tabcolsep}{3pt}
  \begin{tabular}{llll}
    \toprule
    Phase & Epochs & Trained & Learning rate \\
    \midrule
    Probe       & 4       & classifier, Motion Attention & $10^{-3}$ \\
    Fine-tuning & $\le44$ & Res2, Res3, Res5, head       & $5\times10^{-5}$ \\
                &         & classifier, Motion Attention, rings $\Delta$ & $5\times10^{-4}$ \\
    \bottomrule
  \end{tabular}
\end{table}

\textbf{Objective.} For a CutMix batch, a clip $X_a$ receives a box from a
clip $X_b$ covering a fraction $1-\lambda$ of the frame, giving $\tilde{X}$
with targets $y_a$ and $y_b$. With student logits $s = f_\theta(\tilde{X})$
the supervised loss is
\begin{equation}
  \mathcal{L}_{\text{CE}} = \lambda\, \ell(y_a, s) + (1-\lambda)\, \ell(y_b, s),
  \label{eq:ce}
\end{equation}
where $\ell$ is the cross-entropy with label smoothing~\cite{labelsmoothing}
$\varepsilon = 0.1$ ($\lambda = 1$ without CutMix). With teacher logits
$t = f_{\text{T}}(\tilde{X})$ on the same input, the training loss is
\begin{equation}
  \mathcal{L} = (1-\alpha)\, \mathcal{L}_{\text{CE}}
  + \alpha\, \tau^2\, \mathrm{KL}\!\left(\sigma(t/\tau) \,\middle\|\, \sigma(s/\tau)\right),
  \label{eq:kd}
\end{equation}
with $\sigma$ the softmax, $\alpha = 0.5$ and $\tau = 2$; the factor
$\tau^2$ keeps the gradient scale of the second term independent of
$\tau$~\cite{hinton2015}.

\textbf{Teacher.} The teacher is VideoMAE-B~\cite{videomae}, a ViT-B with
86M parameters that takes $16\times224\times224$ clips, pre-trained by
masked autoencoding and fine-tuned on Kinetics-400. For each fold of the
protocol (Section~\ref{sec:protocol}) it is fine-tuned on the training clips
of that fold only: one epoch with the backbone frozen, then all layers at
$3\times10^{-5}$ (classifier $10^{-3}$, weight decay 0.05) for at most 15
epochs, with the same augmentation and checkpoint selection as the student.
It never sees the validation or test clips of the fold. During student
training it runs in inference mode on exactly the clip $\tilde{X}$ that the
student receives, so that its soft target describes the student's input,
including the CutMix box; following Beyer et
al.~\cite{beyer2022}, the student is trained for longer than without
distillation (48 instead of 24 epochs).

\textbf{Optimisation.} AdamW~\cite{adamw} with weight decay $5\times10^{-4}$
and $\beta = (0.9, 0.999)$, a cosine decay within each phase after one
warm-up epoch, gradient clipping at norm 1.0, bf16 arithmetic and batches of
12 clips. Every training clip appears twice per epoch with independent
augmentation. An exponential moving average of the weights (decay 0.999) is
used for validation and as the final model. The checkpoint with the best
validation F1 is kept, and training stops after 16 epochs without
improvement.

\subsection{Inference}

The rings are fused into their kernels and the teacher is discarded. A clip
is classified from one view, the 16 evenly spaced frames. Optionally, $K=4$
views are formed by shifting all frame positions by $k/K$ of the sampling
stride, $k = 0, \dots, K-1$, and their probabilities are averaged. The
threshold is 0.5; Section~\ref{sec:options} shows that tuning it on the
out-of-fold predictions brings no gain.

%======================================================================
\section{Evaluation Protocol}
\label{sec:protocol}

\subsection{Datasets}

We use three public benchmarks (Table~\ref{tab:data}).
RWF-2000~\cite{rwf2000} contains 2{,}000 five-second clips of real fights
and non-violent activity, most of them from surveillance cameras, with an
official split of 1{,}600 training and 400 test clips. Hockey
Fights~\cite{hockey} contains 1{,}000 clips of 41 frames ($360\times288$)
from professional ice-hockey games. RLVS~\cite{rlvs} contains 2{,}000 clips
of street fights and everyday activities collected from YouTube. The last
two have no official split. All clips pass through the pipeline of
Section~\ref{sec:method} and are stored with 64 frames; the Hockey clips,
which are shorter, repeat frames to reach this length.

\subsection{Duplicate audit and splits}

Each clip is described by four evenly spaced frames, converted to grey,
average-pooled to $16\times16$ and standardised; two clips are compared by
the cosine similarity of their descriptors. Pairs above 0.99 are treated as
copies and pairs above 0.9 as the same scene. In RWF-2000, nine training
clips are copies of test clips, all non-violent, and nine further pairs of
copies occur within the training split. We remove these 17 training clips,
so that no test clip and no second copy is used for training. In Hockey
Fights and RLVS we remove 3 and 15 copies. For these two datasets we draw a
test set of 20\% by whole same-scene groups, with both classes balanced,
so that no scene appears on both sides. The copy of RLVS that we obtained
stores each of its 2{,}000 videos twice, in folders named \emph{train} and
\emph{test}; we use one copy. Same-scene pairs that remain in a training
split are kept in the same cross-validation fold.

\begin{table}[t]
  \centering
  \caption{Datasets after the duplicate audit. Groups: same-scene groups
  in the training split (a group never spans two folds). Frames: typical
  length of the original clips (median for RLVS).}
  \label{tab:data}
  \setlength{\tabcolsep}{3pt}
  \begin{tabular}{lcccccc}
    \toprule
    Dataset & Clips & Copies removed & Train & Test & Groups & Frames \\
    \midrule
    RWF-2000      & 2{,}000 & 17 & 1{,}583 & 400 & 1{,}572 & 150 \\
    Hockey Fights & 1{,}000 & 3  & 797     & 200 & 792     & 41 \\
    RLVS          & 2{,}000 & 15 & 1{,}588 & 397 & 1{,}581 & 125 \\
    \bottomrule
  \end{tabular}
\end{table}

\subsection{Cross-validation and decisions}

None of the benchmarks has a validation split. We divide each training
split into four folds that are stratified by class and never separate a
same-scene group; in each fold the checkpoint with the best validation F1
is kept. The four validation sets cover every training clip once, which
gives one out-of-fold (OOF) prediction per clip. Every choice in this
paper, including the architecture, the training recipe, the number of views
and the decision threshold, is made on these OOF predictions. The test
split is scored once. We report the test metrics of single models as the
mean $\pm$ standard deviation over the fold models; each has seen 75\% of
the training clips. Unless stated otherwise, results are averaged over
three training seeds.
% TODO: state the number of seeds per table once the runs are complete.

\subsection{Statistics}

Two models evaluated on the same clips are compared with the exact McNemar
test~\cite{mcnemar} on their paired OOF predictions. We also report the
difference in OOF accuracy with a 95\% bootstrap interval over
clips~\cite{bootstrap}; for several seeds, the per-clip correctness is
averaged over seeds before resampling. On the 1{,}583 OOF clips of
RWF-2000 the half-width of this interval is about 1.1 points in our
comparisons, whereas on its 400 test clips one point corresponds to four
clips.

\subsection{Reproduced baselines}

To compare methods under the same conditions, we train six Kinetics-400
backbones with the same folds, augmentation and recipe as our network
(probe phase, backbone learning rate $5\times10^{-5}$, new layers
$5\times10^{-4}$, 24 epochs, no distillation): X3D-XS, X3D-S and
X3D-M~\cite{x3d} at their native inputs ($4\times160^2$, $13\times160^2$,
$16\times224^2$), and S3D~\cite{s3d}, MC3-18 and
R(2+1)D-18~\cite{r2plus1d} at their training resolutions ($224^2$,
$112^2$, $112^2$) with 16 frames. For the X3D models the stages trained
are those of Table~\ref{tab:phases}; the other backbones are fine-tuned
completely.

\subsection{Metrics, cost and implementation}

We report accuracy, precision, recall and F1 of the violent class, and
ROC-AUC. Cost is given as parameters after fusion, GFLOPs per view counted
as multiply--accumulates with fvcore, latency and throughput on one AMD
MI250 GCD, latency on an AMD EPYC 7413 CPU with 16 threads, and training
time and memory. The code uses PyTorch 2.14 on ROCm 7.2; the channel-wise
3D convolutions run as custom Triton kernels, which make a training step
2.3 times faster than the vendor library with identical results. Each fold
trains on one MI250 GCD.

%======================================================================
\section{Results}
\label{sec:results}

\subsection{Comparison}

Table~\ref{tab:sota} lists published results as reported by their authors
and our models under the protocol of Section~\ref{sec:protocol}. A single
MA-X3D model reaches 87.6\% with one view at 5.46 GFLOPs and 88.1\% with
four views, 1.5--2.0 points above standard fine-tuning of the same network.
Reported lightweight methods reach 89.70--90.25\%. Because their selection
protocol is not stated, and because the published split includes the
duplicated clips described above, the difference should not be read as a
ranking. The VideoMAE-B teacher is 3.6 points more accurate than the
single-view student at 29 times the parameters and 25 times the FLOPs.

\begin{table}[t]
  \centering
  \caption{RWF-2000 test accuracy. Upper part: as reported, selection
  protocol not stated. Lower part: this paper, grouped 4-fold protocol,
  single model (mean $\pm$ std over the fold models).}
  \label{tab:sota}
  \setlength{\tabcolsep}{3pt}
  \begin{tabular}{lccc}
    \toprule
    Method & Params & GFLOPs & Acc. (\%) \\
    \midrule
    ConvLSTM~\cite{convlstm}           & --     & --    & 77.00 \\
    C3D~\cite{c3d}                     & --     & --    & 82.75 \\
    Flow Gated Net~\cite{rwf2000}      & --     & --    & 87.25 \\
    SepConvLSTM~\cite{sepconvlstm}     & --     & --    & 89.75 \\
    MDCN~\cite{mdcn}                   & $<$1M  & --    & 89.70 \\
    BiLTMA~\cite{biltma}               & 4.48M  & --    & 90.25 \\
    VideoMAE~\cite{videomae,mixvit}    & 67--305M & --  & 91.8--92.4 \\
    CUE-Net~\cite{cuenet}              & 354M   & 5826  & 94.00 \\
    \midrule
    X3D-M, final recipe                & 2.98M  & 4.73  & -- \\ % TODO p0_x3d_m
    MA-X3D, standard fine-tuning       & 3.01M  & 5.46  & 86.1 $\pm$ 1.0 \\
    MA-X3D (1 view)                    & 3.01M  & 5.46  & 87.6 $\pm$ 0.7 \\
    MA-X3D (4 views)                   & 3.01M  & 21.9  & 88.1 $\pm$ 0.6 \\
    VideoMAE-B teacher                 & 86.2M  & 135   & 91.2 $\pm$ 0.8 \\
    \bottomrule
  \end{tabular}
  % TODO: fill the missing params/GFLOPs of prior work from their papers.
\end{table}

Table~\ref{tab:main} gives the full metrics. The recipe raises recall on the
violent class from 0.823 to 0.867 at a similar precision, and raises AUC
from 0.923 to 0.949. Averaging the four fold models with four views each
(87.4 GFLOPs per clip) gives 89.25\% accuracy, 0.889 F1 and 0.960 AUC;
since this multiplies the cost by 16, we do not use it as our main result.

\begin{table}[t]
  \centering
  \caption{Out-of-fold (1{,}583 clips) and test metrics (single model,
  mean over the four fold models; std of accuracy in Table~\ref{tab:sota}).
  Violent class for P, R and F1.}
  \label{tab:main}
  \setlength{\tabcolsep}{2.5pt}
  \begin{tabular}{lcccccc}
    \toprule
    & \multicolumn{2}{c}{OOF} & \multicolumn{4}{c}{Test} \\
    \cmidrule(lr){2-3}\cmidrule(lr){4-7}
    Model & Acc & F1 & P & R & F1 & AUC \\
    \midrule
    MA-X3D, standard FT      & 88.44 & 0.885 & 0.892 & 0.823 & 0.856 & 0.923 \\
    MA-X3D (1 view)          & 91.03 & 0.912 & 0.883 & 0.867 & 0.875 & 0.949 \\
    MA-X3D (4 views)         & 91.28 & 0.913 & 0.902 & 0.854 & 0.877 & 0.953 \\
    VideoMAE-B teacher       & 92.48 & 0.926 & 0.902 & 0.925 & 0.913 & 0.967 \\
    \bottomrule
  \end{tabular}
  % TODO: 3-seed values.
\end{table}

\subsection{Where does the improvement come from?}
\label{sec:ablation}

The comparison with conventional fine-tuning mixes two changes, the
network and the training recipe. We separate them with three groups of
experiments, each of which changes one factor. Every comparison is paired:
the two models are trained on the same folds and evaluated on the same
clips. We report the difference in OOF accuracy over the 1{,}583 clips with
a 95\% bootstrap interval, the exact McNemar test on the paired OOF
predictions, and the mean difference of the single fold models on the test
split. Rows with three seeds report the mean over the 12 fold models.
% TODO: 3-seed values for every row marked [s0].

\textbf{Network and recipe (Table~\ref{tab:factorial}).} X3D-M and MA-X3D
are each trained with the conventional recipe and with the final recipe.
If the modules only helped because of the recipe, or the recipe only
because of the modules, the gains would not appear in both columns.
% TODO: fill Table IV from q0 (X3D-M conventional), e00, p0, e11 (3 seeds each) and
% describe: gain of the modules under each recipe, gain of the recipe for each network.

\begin{table}[t]
  \centering
  \caption{Network $\times$ training recipe (1 view). OOF accuracy over
  1{,}583 clips / single-model test accuracy, mean over seeds and folds.}
  \label{tab:factorial}
  \setlength{\tabcolsep}{4pt}
  \begin{tabular}{lcc}
    \toprule
    & Conventional recipe & Final recipe \\
    \midrule
    X3D-M   & -- / --          & -- / -- \\ % TODO q0 | p0
    MA-X3D  & 88.44 / 86.1     & 91.03 / 87.6 \\ % [s0] e00 | e11
    \bottomrule
  \end{tabular}
\end{table}

\textbf{Components and controls (Table~\ref{tab:ablation}).} Under the
final recipe we add each module to X3D-M alone, and we compare each module
with a control that has the same parameters and cost but lacks the
property we attribute the gain to. For Motion Attention, the control keeps
the gate but feeds it an all-zero motion map, so any gain of the control
comes from the extra gating parameters rather than from motion. For the
wide kernel, the first control inflates the kernel to a dense $5\times5$
tensor with a zero ring, so the ring shares the backbone learning rate; the
second keeps the reparameterisation but places the ring in the backbone
group. A third check needs no training: we set the learned ring of the
full model to zero at test time. This lowers single-model accuracy from
87.6\% to 80.6\%; the ring carries 30\% of the kernel magnitude (mean over
the eight widened convolutions), so the trained network relies on the added
spatial support.
% TODO: describe p0, p1, p2, p6, p4, p5 with differences, intervals and p-values.

\begin{table}[t]
  \centering
  \caption{Components and controls under the final recipe (1 view).
  $\Delta$: OOF accuracy difference to X3D-M with 95\% bootstrap interval;
  $p$: exact McNemar test against X3D-M on the OOF predictions.}
  \label{tab:ablation}
  \setlength{\tabcolsep}{2pt}
  \begin{tabular}{lccccc}
    \toprule
    Model & Params & OOF Acc & $\Delta$ & $p$ & Test Acc \\
    \midrule
    X3D-M                              & 2.98M & -- & -- & -- & -- \\ % TODO p0
    + MA                               & 2.98M & -- & -- & -- & -- \\ % TODO p1
    + MA, motion map set to zero       & 2.98M & -- & -- & -- & -- \\ % TODO p6
    + WK                               & 3.01M & -- & -- & -- & -- \\ % TODO p2
    + WK, dense $5\times5$             & 3.01M & -- & -- & -- & -- \\ % TODO p4
    + WK, ring at backbone rate        & 3.01M & -- & -- & -- & -- \\ % TODO p5
    + MA + WK (MA-X3D)                 & 3.01M & 91.03 & -- & -- & 87.6 \\ % [s0]
    \quad ring set to zero at test     & 3.01M & -- & -- & -- & 80.6 \\
    \bottomrule
  \end{tabular}
\end{table}

\textbf{Training recipe (Table~\ref{tab:recipe}).} For the same network we
change the recipe one step at a time. Raising the backbone learning rate
from $10^{-5}$ to $5\times10^{-5}$ adds 1.26 OOF points ($[+0.13, +2.40]$,
$p=0.040$) and 1.6 points on the test split: with the lower rate, training
accuracy stayed below validation accuracy and most folds stopped early,
which suggests that the backbone was not adapting enough. Distillation with
the conventional 24-epoch schedule leaves OOF accuracy unchanged. With 48
epochs it adds 1.33 OOF points over R1 ($[+0.32, +2.34]$, $p=0.015$), and
the selected epochs move from 9--20 to 20--45. Because R3 differs from R1
in both the teacher and the schedule length, the control R1$'$ trains for
48 epochs without the teacher; the effect of distillation is the
difference between R3 and R1$'$. On the test split, the single-model
difference between R3 and R1 is $-0.06 \pm 0.67$ points, so the OOF gain of
distillation is not visible there at this sample size.
% TODO: e19 (R1') and 3 seeds; rewrite the last two sentences accordingly.

The lower part of Table~\ref{tab:recipe} lists changes that were tried and
not adopted because they did not improve OOF accuracy over the final recipe
or cost more: a larger distillation weight, a higher backbone rate,
fine-tuning Res4 ($+0.82$ OOF points over R1, $p=0.17$), the deeper X3D-L
backbone (5.4M parameters), and 32 frames per clip (10.9 GFLOPs).

\begin{table}[t]
  \centering
  \caption{Training recipe for the same network (MA-X3D, 1 view, threshold
  0.5). $p$: McNemar test against the row above on the OOF predictions.
  Test: single model, mean $\pm$ std over the fold models.}
  \label{tab:recipe}
  \setlength{\tabcolsep}{2pt}
  \begin{tabular}{llccc}
    \toprule
    & Recipe & OOF Acc & $p$ & Test Acc \\
    \midrule
    R0 & conventional (lr $10^{-5}$, 24 ep.)        & 88.44 & --    & 86.1 $\pm$ 1.0 \\
    R1 & backbone lr $5\times10^{-5}$               & 89.70 & 0.040 & 87.7 $\pm$ 0.9 \\
    R1$'$ & R1, 48 ep., no teacher                  & --    & --    & -- \\ % TODO e19
    R2 & R1 + distillation, 24 ep.                  & 89.83 & --    & 87.2 $\pm$ 0.8 \\
    R3 & R1 + distillation, 48 ep. (final)          & 91.03 & 0.015 & 87.6 $\pm$ 0.7 \\
    \midrule
    & R2 with $\alpha = 0.9$                        & 90.27 & --    & 88.1 $\pm$ 0.5 \\
    & R1 with lr $10^{-4}$                          & 89.89 & --    & 87.4 $\pm$ 1.1 \\
    & R1 + Res4 fine-tuned                          & 90.52 & 0.17  & 88.2 $\pm$ 0.7 \\
    & R1 with X3D-L backbone                        & 89.83 & --    & 87.2 $\pm$ 0.5 \\
    & R0 with 32 frames                             & 88.76 & --    & 85.8 $\pm$ 2.6 \\
    \bottomrule
  \end{tabular}
  % p for R3 is against R1 (R2 is shown for reference). TODO: 3 seeds;
  % add e16 (96 ep.), e17 (48 ep., alpha 0.9), e18 (48 ep. + Res4) if they change a decision.
\end{table}

\subsection{Motion Attention design}

% TODO (seed 0, final recipe): one step only (p7), raw map without temporal modes
% (p8), gate after Res3 (p9) and after Res5 (p10), each against MA-X3D with the
% McNemar test. Keep as a short table or two sentences, depending on space.

\subsection{Inference options}
\label{sec:options}

On the OOF predictions, four views raise accuracy by 0.25 points over one
view (0.4 points on the test split) at four times the cost. Horizontal-flip
test-time augmentation lowers OOF accuracy slightly, and tuning the
threshold on $[0.30, 0.70]$ selects 0.50--0.52 with no measurable gain. We
therefore use one view and a threshold of 0.5 for deployment and report
four views as an option.

\subsection{Qualitative analysis}

% TODO: Grad-CAM with and without Motion Attention, gate maps on violent and
% non-violent clips, and typical errors (crowded scenes, distant actors,
% sport and horseplay). Generate with `ma-x3d visualize` on the final model.

\subsection{Cost}

Table~\ref{tab:cost} compares the deployed student with the teacher. The
student has 29 times fewer parameters and 25 times fewer FLOPs, and is
about twice as fast on CPU. On the MI250 GPU the two models have similar
latency and throughput: depthwise 3D convolutions are limited by memory
bandwidth on this hardware, while the dense matrix products of the
transformer run close to peak. The advantage of the student is therefore
largest on CPUs and memory-limited devices.
% TODO: one edge-device measurement (e.g. Jetson or laptop CPU).

\begin{table}[t]
  \centering
  \caption{Cost of one model at $16\times224\times224$. GPU: one AMD MI250
  GCD shared with other jobs; CPU: AMD EPYC 7413, 16 threads, PyTorch.}
  \label{tab:cost}
  \setlength{\tabcolsep}{3pt}
  \begin{tabular}{lcc}
    \toprule
    & MA-X3D & VideoMAE-B \\
    \midrule
    Parameters                         & 3.01M & 86.2M \\
    GFLOPs per view                    & 5.46  & 135 \\
    GPU latency, batch 1, fp32 (ms)    & 30    & 24 \\
    GPU throughput, batch 16, fp16 (clips/s) & 118 & 136 \\
    CPU latency, batch 1 (ms)          & 208   & 400 \\
    Training per fold (min)            & 48    & 12 \\
    Peak training memory (GB)          & 5.8   & 9.4 \\
    \bottomrule
  \end{tabular}
\end{table}

%======================================================================
\section{Limitations}

The results come from one benchmark whose test split has 400 clips, so one
point of accuracy is four clips. For every model, including the teacher,
test accuracy is 1--3 points below OOF accuracy, which suggests that the
test clips are harder than the training clips. The static crop assumes a
fixed camera. Training requires a teacher that is 29 times larger than the
student. The model classifies trimmed clips and does not localise events in
time. The Hockey Fight results of our earlier work and the effect of the
person-centred crop were not re-evaluated under the present protocol.

%======================================================================
\section{Conclusion}
\label{sec:conclusion}

MA-X3D adds a frame-difference gate and a fused wide kernel to X3D-M at a
cost of 36K parameters. Trained with a higher backbone learning rate and
long, view-consistent distillation from a video transformer, a single model
with 3.0M parameters reaches 87.6\% on RWF-2000 with one view, under a
protocol in which the test split takes no part in any decision. Future work
includes temporal localisation on untrimmed streams, measurements on edge
devices, and cross-dataset evaluation under the same protocol.

\bibliographystyle{IEEEtran}
\bibliography{refs}

\end{document}
